% Options for packages loaded elsewhere
% Options for packages loaded elsewhere
\PassOptionsToPackage{unicode}{hyperref}
\PassOptionsToPackage{hyphens}{url}
\PassOptionsToPackage{dvipsnames,svgnames,x11names}{xcolor}
%
\documentclass[
  letterpaper,
  DIV=11,
  numbers=noendperiod]{scrartcl}
\usepackage{xcolor}
\usepackage{amsmath,amssymb}
\setcounter{secnumdepth}{-\maxdimen} % remove section numbering
\usepackage{iftex}
\ifPDFTeX
  \usepackage[T1]{fontenc}
  \usepackage[utf8]{inputenc}
  \usepackage{textcomp} % provide euro and other symbols
\else % if luatex or xetex
  \usepackage{unicode-math} % this also loads fontspec
  \defaultfontfeatures{Scale=MatchLowercase}
  \defaultfontfeatures[\rmfamily]{Ligatures=TeX,Scale=1}
\fi
\usepackage{lmodern}
\ifPDFTeX\else
  % xetex/luatex font selection
\fi
% Use upquote if available, for straight quotes in verbatim environments
\IfFileExists{upquote.sty}{\usepackage{upquote}}{}
\IfFileExists{microtype.sty}{% use microtype if available
  \usepackage[]{microtype}
  \UseMicrotypeSet[protrusion]{basicmath} % disable protrusion for tt fonts
}{}
\makeatletter
\@ifundefined{KOMAClassName}{% if non-KOMA class
  \IfFileExists{parskip.sty}{%
    \usepackage{parskip}
  }{% else
    \setlength{\parindent}{0pt}
    \setlength{\parskip}{6pt plus 2pt minus 1pt}}
}{% if KOMA class
  \KOMAoptions{parskip=half}}
\makeatother
% Make \paragraph and \subparagraph free-standing
\makeatletter
\ifx\paragraph\undefined\else
  \let\oldparagraph\paragraph
  \renewcommand{\paragraph}{
    \@ifstar
      \xxxParagraphStar
      \xxxParagraphNoStar
  }
  \newcommand{\xxxParagraphStar}[1]{\oldparagraph*{#1}\mbox{}}
  \newcommand{\xxxParagraphNoStar}[1]{\oldparagraph{#1}\mbox{}}
\fi
\ifx\subparagraph\undefined\else
  \let\oldsubparagraph\subparagraph
  \renewcommand{\subparagraph}{
    \@ifstar
      \xxxSubParagraphStar
      \xxxSubParagraphNoStar
  }
  \newcommand{\xxxSubParagraphStar}[1]{\oldsubparagraph*{#1}\mbox{}}
  \newcommand{\xxxSubParagraphNoStar}[1]{\oldsubparagraph{#1}\mbox{}}
\fi
\makeatother


\usepackage{longtable,booktabs,array}
\usepackage{calc} % for calculating minipage widths
% Correct order of tables after \paragraph or \subparagraph
\usepackage{etoolbox}
\makeatletter
\patchcmd\longtable{\par}{\if@noskipsec\mbox{}\fi\par}{}{}
\makeatother
% Allow footnotes in longtable head/foot
\IfFileExists{footnotehyper.sty}{\usepackage{footnotehyper}}{\usepackage{footnote}}
\makesavenoteenv{longtable}
\usepackage{graphicx}
\makeatletter
\newsavebox\pandoc@box
\newcommand*\pandocbounded[1]{% scales image to fit in text height/width
  \sbox\pandoc@box{#1}%
  \Gscale@div\@tempa{\textheight}{\dimexpr\ht\pandoc@box+\dp\pandoc@box\relax}%
  \Gscale@div\@tempb{\linewidth}{\wd\pandoc@box}%
  \ifdim\@tempb\p@<\@tempa\p@\let\@tempa\@tempb\fi% select the smaller of both
  \ifdim\@tempa\p@<\p@\scalebox{\@tempa}{\usebox\pandoc@box}%
  \else\usebox{\pandoc@box}%
  \fi%
}
% Set default figure placement to htbp
\def\fps@figure{htbp}
\makeatother


% definitions for citeproc citations
\NewDocumentCommand\citeproctext{}{}
\NewDocumentCommand\citeproc{mm}{%
  \begingroup\def\citeproctext{#2}\cite{#1}\endgroup}
\makeatletter
 % allow citations to break across lines
 \let\@cite@ofmt\@firstofone
 % avoid brackets around text for \cite:
 \def\@biblabel#1{}
 \def\@cite#1#2{{#1\if@tempswa , #2\fi}}
\makeatother
\newlength{\cslhangindent}
\setlength{\cslhangindent}{1.5em}
\newlength{\csllabelwidth}
\setlength{\csllabelwidth}{3em}
\newenvironment{CSLReferences}[2] % #1 hanging-indent, #2 entry-spacing
 {\begin{list}{}{%
  \setlength{\itemindent}{0pt}
  \setlength{\leftmargin}{0pt}
  \setlength{\parsep}{0pt}
  % turn on hanging indent if param 1 is 1
  \ifodd #1
   \setlength{\leftmargin}{\cslhangindent}
   \setlength{\itemindent}{-1\cslhangindent}
  \fi
  % set entry spacing
  \setlength{\itemsep}{#2\baselineskip}}}
 {\end{list}}
\usepackage{calc}
\newcommand{\CSLBlock}[1]{\hfill\break\parbox[t]{\linewidth}{\strut\ignorespaces#1\strut}}
\newcommand{\CSLLeftMargin}[1]{\parbox[t]{\csllabelwidth}{\strut#1\strut}}
\newcommand{\CSLRightInline}[1]{\parbox[t]{\linewidth - \csllabelwidth}{\strut#1\strut}}
\newcommand{\CSLIndent}[1]{\hspace{\cslhangindent}#1}



\setlength{\emergencystretch}{3em} % prevent overfull lines

\providecommand{\tightlist}{%
  \setlength{\itemsep}{0pt}\setlength{\parskip}{0pt}}



 


\usepackage{hyperref, mathtools, amsmath, scrlayer-scrpage, makeidx} \rohead{ } \lofoot{ \vspace*{-150px} \centering \tiny{ \vspace*{3px} \copyright \ \ Update at \href{https://eugenionet.github.io/physecon/}{eugenionet.github.io/physecon}, eugenionetweb@gmail.com } } \makeindex
\KOMAoption{captions}{tableheading}
\makeatletter
\@ifpackageloaded{caption}{}{\usepackage{caption}}
\AtBeginDocument{%
\ifdefined\contentsname
  \renewcommand*\contentsname{Table of contents}
\else
  \newcommand\contentsname{Table of contents}
\fi
\ifdefined\listfigurename
  \renewcommand*\listfigurename{List of Figures}
\else
  \newcommand\listfigurename{List of Figures}
\fi
\ifdefined\listtablename
  \renewcommand*\listtablename{List of Tables}
\else
  \newcommand\listtablename{List of Tables}
\fi
\ifdefined\figurename
  \renewcommand*\figurename{Figure}
\else
  \newcommand\figurename{Figure}
\fi
\ifdefined\tablename
  \renewcommand*\tablename{Table}
\else
  \newcommand\tablename{Table}
\fi
}
\@ifpackageloaded{float}{}{\usepackage{float}}
\floatstyle{ruled}
\@ifundefined{c@chapter}{\newfloat{codelisting}{h}{lop}}{\newfloat{codelisting}{h}{lop}[chapter]}
\floatname{codelisting}{Listing}
\newcommand*\listoflistings{\listof{codelisting}{List of Listings}}
\makeatother
\makeatletter
\makeatother
\makeatletter
\@ifpackageloaded{caption}{}{\usepackage{caption}}
\@ifpackageloaded{subcaption}{}{\usepackage{subcaption}}
\makeatother
\usepackage{bookmark}
\IfFileExists{xurl.sty}{\usepackage{xurl}}{} % add URL line breaks if available
\urlstyle{same}
\hypersetup{
  pdftitle={Physics Concepts as Indicative Framework to Model Systems of Economies Ex-ante},
  pdfauthor={Eugenio Net},
  pdfkeywords={Physics, Economy},
  colorlinks=true,
  linkcolor={blue},
  filecolor={Maroon},
  citecolor={Blue},
  urlcolor={Blue},
  pdfcreator={LaTeX via pandoc}}


\title{Physics Concepts as Indicative Framework to Model Systems of
Economies Ex-ante}
\usepackage{etoolbox}
\makeatletter
\providecommand{\subtitle}[1]{% add subtitle to \maketitle
  \apptocmd{\@title}{\par {\large #1 \par}}{}{}
}
\makeatother
\subtitle{Coincise Notes and Conjectures Explorations of Natural Laws as
Approach to Preliminarly Structure Economies}
\author{Eugenio Net}
\date{2026-07-02}
\begin{document}
\maketitle
\begin{abstract}
An exploratory attempt \ldots{} Interconnection network between
movement, force, momentum, energy, and work. Describing systems with
Lagrangian and Hamltonian when constraints and starting conditions
exist. Modeling systems and optimisation with constraints and
conditions. Interwoven system of macro- and microeconomie under
constraints, limitations, and conditions.
\end{abstract}

\renewcommand*\contentsname{Table of contents}
{
\hypersetup{linkcolor=}
\setcounter{tocdepth}{3}
\tableofcontents
}

\begin{center}\rule{0.5\linewidth}{0.5pt}\end{center}

\hfill\break
\hfill\break

\emph{WORKING NOTES DRAFT}

\hfill\break
\hfill\break

\section{Introduction}\label{sec-introduction}

Nature strives through Physics to behave economical by using its
resources in an austere way: Principle of Least Action.

Economy - not pre-given by natural laws - is constructed as social set
of co-living rules to allocate, distribute and share (natural and
non-natural) resources. Ex-ante modelling economic systems by rules
(contractual, legislative, judicative, executive).

The state (situation, position) of the Economy can be described by
e.g.~Assets, Liabilities, Expenses, Incomes, Profitability,
Productivity, etc.. The Economy changes through time from a state
\(r_0=\) ``start'' to a state \(r_1=\) ``end''. For the state to change,
work is applied by processing limited resources (e.g.~access to raw
materials and availability of intermediate goods), constrained capital
(e.g.~financial funds, interest rates, and inflation), as well as labor
capacity (e.g machines and human force).

How can a trajectory \(r = r_1 - r_0\) between two states be modelled by
work \(W\), which is applied for the change of the state to occure:
searching the condition by defing target state through ``path of chage''
\(=r\)?

\begin{equation}\phantomsection\label{eq-intro}{\begin{split} &
r = f(x_i) = \int A = \int \int W = \int \int (\mathcal{L} + \mathcal{H}) = \int \int (\Delta E - \Delta H)
\end{split}}\end{equation}

\hfill\break
\hfill\break

\section{Motion and State Change}\label{sec-motion}

\begin{itemize}
\tightlist
\item
  motion of matter in space and time
\item
  change of state
\item
  displacement
\item
  diffusion
\item
  dissipation ad refraction
\item
  distance, velocity, momentum, acceleration (curvature), force, energy,
  work, and action
\item
  isotropy and homomenity
\end{itemize}

Motion:
\begin{equation}\phantomsection\label{eq-trajectory}{\begin{split} &
r = r(t) = r_x(t_{x}) - r_0(t_{0} ) = \Delta r(\Delta t)
\\& \stackrel{Path}{=}
State(arrival) - State(start) = Trajectory(Change)
\end{split}}\end{equation}

\begin{equation}\phantomsection\label{eq-motion}{\begin{split} &
v = \dot r = \frac{d}{dt}r
\qquad
p = \frac{1}{v} \sqrt{E^2 - (mv^2)^2}
\qquad
F = \dot p = \frac{d}{dt}p
\end{split}}\end{equation}

Energy total:
\begin{equation}\phantomsection\label{eq-energytotal}{\begin{split} &
E = Motion_{Energy} + Rest_{Energy} = E_{kin} + E_{pot}
\end{split}}\end{equation}

\hfill\break
\hfill\break

\section{Work and Energy}\label{sec-workenergy}

\(\exists\) Work done or received \(W \neq 0 \Rrightarrow \exists\)
Change of State (``Displacement'', ``Deformation'') \(r \neq 0\):

\begin{equation}\phantomsection\label{eq-energy}{\begin{split} &
E = \mathcal{K} + \mathcal{U}  \qquad  E_{kin} \coloneqq \mathcal{K} = \frac{1}{2}pv   \qquad   E_{pot} \coloneqq \mathcal{U} = \mathcal{V} + \mathcal{Z}
\end{split}}\end{equation}

Energy is conserved in konservative systems; There is no work in closed
path between two points (loop) of any shape under path variation, hence
the ``Principle of Action'' applies (extremities of bumps and saddles
exist): \(\delta A = 0\)

Potentials:
\begin{equation}\phantomsection\label{eq-potential}{\begin{split} &
\mathcal{U} = \mathcal{V} + \mathcal{Z}
\\&
\mathcal{V} 
    \begin{cases}
      = \frac{1}{2} kr^2 : \text{isotrop harmonical oscilator central field Potential, } \\
        \qquad \text{closed bounded path, } \text{"spring constant" } k = m\omega^2 = 2\pi \cdot \nu  \\
      = -\frac{1}{r}\alpha: \text{gravitational central field Potential, } \\
        \qquad \text{closed bounded path, } \alpha = \left( \frac{2\pi}{\Delta \varphi} \right)^2 = const. \\
      \sim \frac{1}{r}e^{\mu r} : \text{Yukawa Potential}, \text{"Photon mass" } m_{\gamma} = \mu \frac{\hbar}{c} \\
      ...
    \end{cases}
\end{split}}\end{equation}

Centrifugal Potential:
\begin{equation}\phantomsection\label{eq-centrifugalpot}{\begin{split} &
\mathcal{Z} = \frac{1}{2} \frac{1}{m} \left( \frac{L_0}{r} \right)^2 \text{, angular momentum}=L_0
\end{split}}\end{equation}

Work:
\begin{equation}\phantomsection\label{eq-workaction}{\begin{split} &
W = \partial_t A = \partial_t (E \cdot t) = \partial_t (p ⋅ r) = F_1 \cdot r_1 + p_2 \cdot v_2 = 
    \begin{cases}
      F \cdot r \\
      p \cdot v
    \end{cases}
\end{split}}\end{equation}

\emph{Geometrical}:

\begin{equation}\phantomsection\label{eq-workgeom}{\begin{split} &
W = F \cdot r = mr \cdot a = I ⋅ \frac{v}{t} =  \frac{p}{t} ⋅ tv = Ft \cdot v =  pv = mv^2 
\end{split}}\end{equation}

\emph{Analytical}:

\begin{equation}\phantomsection\label{eq-lagrange}{\begin{split} &
\mathcal{L} = \mathcal{L}(r,\dot r,t) = E_{kin} - E_{pot} 
\end{split}}\end{equation}

\begin{equation}\phantomsection\label{eq-hamilton}{\begin{split} &
\mathcal{H} = \mathcal{H}(r,p,t) = \sum_i^f \dot r_i p_i - \mathcal{L}
\stackrel{konserv.}{=}
E_{kin} + E_{pot} = const.
\end{split}}\end{equation}

\begin{equation}\phantomsection\label{eq-action}{\begin{split} &
A = E \cdot t = p \cdot r
\stackrel{\mathcal{H}=0}{=}
\int_{t_0}^{t_x} \mathcal{L} dt
\end{split}}\end{equation}

\begin{equation}\phantomsection\label{eq-workanal}{\begin{split} &
W = p \cdot v = \mathcal{L} + \mathcal{H} = (\mathcal{K} - \mathcal{U}) + \mathcal{H}
\\& \stackrel{konserv.}{=}
(\mathcal{K} - \mathcal{U}) + (\mathcal{K} + \mathcal{U}) = (\mathcal{K} - \mathcal{U}) + E
\end{split}}\end{equation}

\emph{Electro Magnetic}:

\begin{equation}\phantomsection\label{eq-workelmag}{\begin{split} &
W = \mathcal{L} + \mathcal{H} = \left( \frac{1}{2}mv^2 - q(\Phi - \vec{\mathcal{A}} \cdot \vec{v}) \right) + \left( \frac{1}{2} \frac{1}{m} (\vec{p} - q\vec{\mathcal{A}})^2 + q\Phi \right)\end{split}}\end{equation}

\emph{Relativistic}:

\begin{equation}\phantomsection\label{eq-workrel}{\begin{split} &
W = p \cdot v = \frac{v}{c} \sqrt{E^2 - (mc^2)^2}
\end{split}}\end{equation}

\emph{Thermal}:

\begin{equation}\phantomsection\label{eq-workthermal}{\begin{split} &
W = Fr = PV^{\iota} = T \cdot V^{\frac{f+2}{f}} = T \cdot V^{(\frac{c_P}{c_V}-1)} = TV^{\kappa-1}  \\&
= T \cdot l \cdot g_0 = T \cdot n \cdot b_0 = T \cdot m \cdot c_0
\end{split}}\end{equation}

\emph{Thermo Statistical}:

\begin{equation}\phantomsection\label{eq-workthemostat}{\begin{split} &
W = \mathcal{L} + \mathcal{H} = ΔU - ΔH = \Delta(T \cdot \Delta S - P \cdot \Delta V) - \Delta(T\cdot S)
\end{split}}\end{equation}

\emph{Quantistic}:

\begin{equation}\phantomsection\label{eq-workquant}{\begin{split} &
W = p \cdot v = ħ\mathcal{k} \cdot ω = ℎν\mathcal{k} = E\mathcal{k}
\end{split}}\end{equation}

\begin{equation}\phantomsection\label{eq-uncertainity}{\begin{split} &
\Delta A = \Delta p  \cdot \Delta r = \Delta E \cdot \Delta t \ge \frac{1}{2} \hbar
\end{split}}\end{equation}

\hfill\break
\hfill\break

see description of ``movement'' Equation~\ref{eq-movement}

acronyms see Section~\ref{sec-acroworkenergy}

\hfill\break
\hfill\break

\section{Constrains and Conditions}\label{sec-constraints}

Constraints:

\begin{itemize}
\tightlist
\item
  initial state, starting conditions
\item
  acting boundary conditions (to move within): e.g.~demography,
  economical and political stability, labor force, education and
  know-how, health and technological level
\item
  access to limited resources: e.g.~chemical elements, minerals, fuels,
  commodities, trade and transportations, etc.
\item
  contractual agreements and covenants, legal assurance
\item
  juristic law framework
\end{itemize}

Modelling movement (change of state) under constraints:

\begin{enumerate}
\def\labelenumi{\arabic{enumi}.}
\item
  Movement ist charecterized by change of energy, Lagrangian
  \(\mathcal{L} \neq 0\):
  \begin{equation}\phantomsection\label{eq-lagrangian}{\begin{split} &
  \mathcal{L} = \mathcal{K} - \mathcal{U}
  \qquad
  \frac{\partial \mathcal{L}}{\partial r} = F
  \qquad
  \frac{\partial \mathcal{L}}{\partial v} = p
  \end{split}}\end{equation}
\item
  Movement happens within constraints of forces. Constrains are
  observable, whereas forces might be unknown. Therefore, \(r\) can be
  determined from \(\mathcal{L}\) (Energy and Work) of the
  \textbf{system at rest \(F=0\)}, instead than from ``searching'' all
  acting forces \(F_i\). In general, unknow forces at motion are defined
  by solving the Eigenvalueproblem. Solving for \textbf{\(\mathcal{L}\)}
  at equilibirium:
  \begin{equation}\phantomsection\label{eq-constrainforce}{\begin{split} &
  F_{total} = \hat{F_i} + \tilde{F_i} = \hat{F_i} + \sum_{k=1}^{s_d} \lambda_k a_{ki} = \frac{d}{dt} \left( \frac{\partial \mathcal{L}}{\partial v_i} \right) - \frac{\partial \mathcal{L}}{\partial r_i} 
  = \begin{cases}
   = 0    : \text{ Equilibrium} \\
   \neq 0 : \text{ Movement}
    \end{cases}
  \end{split}}\end{equation}

  and at same time solving for constraining conditions that hold the
  system at rest (equilibrium):
  \begin{equation}\phantomsection\label{eq-constraincondition}{\begin{split} &
  \sum_{l=1}^f a_{kl} v_l + b_k = 0
  \end{split}}\end{equation}
\item
  Rest is characterized by unchanged energy, Hamiltonian
  \(\mathcal{H} = E_{kin} + E_{pot} = const.\). There is movement from
  state of rest if the new \(\mathcal{H} \neq 0\) differs from the
  initial one. Using the found solution for the Lagrangian
  \(\mathcal{L}\) containing the constraints, the Hamiltonian for the
  energy is deduced:
  \begin{equation}\phantomsection\label{eq-hamiltonian}{\begin{split} &
  \mathcal{H} = \sum_i p_i v_i - \mathcal{L} = pv - \mathcal{L}  = Fr - \mathcal{L} = W - \mathcal{L}  = (\Delta E - \Delta H) - \mathcal{L}
  \end{split}}\end{equation}
\item
  Change of state/position is characterized by movement. Eventually, by
  solving the new system of differential equations, from the found
  solution of the energy \(\mathcal{H}\), the set \({r,v,p}\) which
  describe the movement (and thus the change of state), is deduced:
  \begin{equation}\phantomsection\label{eq-movement}{\begin{split} &
  \dot r_i = v_i = \frac{\partial\mathcal{H}}{\partial p_i}   \qquad
  \dot p_i = F_i = - \frac{\partial\mathcal{H}}{\partial r_i}
  \end{split}}\end{equation}
\end{enumerate}

see description of ``work'' Equation~\ref{eq-workgeom}

where:

\(s_d =\) constraints   \(i = 1,...,f =\) dedrees of freedom  
\(k = 1,...,s_d =\) constraints

\(\hat{F_i} =\) known constraining forces  
\(\tilde{F_i} = \sum_{k=1}^{s_d} \lambda_k a_{ki} =\) unknown
constraining forces

\(\lambda_k(t) =\) Eigenvalues for transformation of reference system

\(a_{ki} =\) Eigenvectors of transformation matrix

\(\lambda_k \cdot \vec r_i = \lambda_k \cdot \left[ a_{ki} \right]_{matrix}\)

\hfill\break
\hfill\break

\section{Determinism and Uncertainity}\label{sec-determinism}

\begin{itemize}
\tightlist
\item
  Determinism through contratcting \& jurisdiction
\item
  Uncertainity minimizsations
\item
  Statistical approximation
\item
  Prognosis matching Determinism
\end{itemize}

see Equation~\ref{eq-uncertainity}

\hfill\break
\hfill\break

\section{Ex-ante Framing and Post-ante
Optimization}\label{sec-optimization}

Framing contracts and jurisdiction ahead to match
\(p = \text{"Economy Funds"}\) and \(v = \text{"Economy Velocity"}\)
that can be modelled with typical approaches for fuction
\(z = \text{"arrival state"}\) depending on
\(x = \text{"starting state"}\), so that ``work''
\(W = F \cdot r =  F \cdot (z - x)\):

\begin{longtable}[]{@{}
  >{\raggedright\arraybackslash}p{(\linewidth - 4\tabcolsep) * \real{0.2329}}
  >{\raggedright\arraybackslash}p{(\linewidth - 4\tabcolsep) * \real{0.2740}}
  >{\raggedright\arraybackslash}p{(\linewidth - 4\tabcolsep) * \real{0.4932}}@{}}
\caption{Modelling Functions}\label{tbl-modeling}\tabularnewline
\toprule\noalign{}
\begin{minipage}[b]{\linewidth}\raggedright
Type
\end{minipage} & \begin{minipage}[b]{\linewidth}\raggedright
Model
\end{minipage} & \begin{minipage}[b]{\linewidth}\raggedright
Linearisation
\end{minipage} \\
\midrule\noalign{}
\endfirsthead
\toprule\noalign{}
\begin{minipage}[b]{\linewidth}\raggedright
Type
\end{minipage} & \begin{minipage}[b]{\linewidth}\raggedright
Model
\end{minipage} & \begin{minipage}[b]{\linewidth}\raggedright
Linearisation
\end{minipage} \\
\midrule\noalign{}
\endhead
\bottomrule\noalign{}
\endlastfoot
Linear & \(z = a + bx\) & \(z = a + bx\) \\
Gemetric & \(z = a \cdot x^b\) &
\(\log z = \log a + b \cdot (\log x)\) \\
Exponential deca. & \(z = a \cdot b^x\) &
\(\log z = \log a + (\log b) \cdot x\) \\
Exponential nat. & \(z = a \cdot e^{bx}\) & \(\ln z = \ln a + bx\) \\
Periodical & & \\
Polynomial & & \\
Taylor & & \\
\end{longtable}

\hfill\break
\hfill\break
Linear optimization to achieve target states ``\(z, h\)'':

\begin{enumerate}
\def\labelenumi{\arabic{enumi}.}
\item
  Oneness situaiton modelling:

  minimize efforts = \(z\) or maximizie return = \(z\)
\end{enumerate}

\begin{equation}\phantomsection\label{eq-situationopt}{\begin{split} &
\text{target:}            \hphantom{min.nstraints:}     z = c^T \cdot x \\&
\text{min. constraints:}  \hphantom{max.}               \left[ a_{ki} \right] \cdot x \ge b   \\&
\text{max. constraints:}  \hphantom{min.}               \left[ a_{ki} \right] \cdot x \le b   \\&
\text{condition:}         \hphantom{max.raints:}        x \ge 0   \\&
\end{split}}\end{equation}

\begin{enumerate}
\def\labelenumi{\arabic{enumi}.}
\setcounter{enumi}{1}
\item
  Duality situation modelling:

  minimize efforts = \(z\) and maximize return = \(h\)
\end{enumerate}

\begin{equation}\phantomsection\label{eq-situationdual}{\begin{split} &
\text{target min.:}         \hphantom{ andints:}      z = c^T \cdot x     \\&
\text{min. constraints:}    \hphantom{ and}           \left[ a_{ki} \right] \cdot x \ge b   \\&
\text{condition:}           \hphantom{ andraints:}    x \ge 0       \\&
\textit{ and}    \\&
\text{target max.:}         \hphantom{ andints:}      h = b^T \cdot y \\&
\text{max. constraints:}    \hphantom{ and}           \left[ a_{ki} \right]^T \cdot y \le c   \\&
\text{condition:}           \hphantom{ andraints:}    y \ge 0         \\&
\end{split}}\end{equation}

\begin{enumerate}
\def\labelenumi{\arabic{enumi}.}
\setcounter{enumi}{2}
\tightlist
\item
  Normalization:
\end{enumerate}

\begin{equation}\phantomsection\label{eq-optnorm}{\begin{split} &
\text{target:}              \hphantom{min.nstraints:}      z = c^T \cdot x \\&
\text{min. constraints:}    \hphantom{max.}    \left( \left[ a_{ki} \right] - x_{n+k} \right) \cdot x = b   \\&
\text{max. constraints:}    \hphantom{min.}    \left( \left[ a_{ki} \right] + x_{n+k} \right) \cdot x = b   \\&
\text{condition:}           \hphantom{max.raints:}    x \ge 0       \\&
\end{split}}\end{equation}

\begin{enumerate}
\def\labelenumi{\arabic{enumi}.}
\setcounter{enumi}{3}
\tightlist
\item
  Solution with
  \href{https://cbom.atozmath.com/CBOM/Simplex.aspx?q=sm}{Simplex
  algorithm}
\end{enumerate}

multilinear and exponential models \ldots{}

\hfill\break
\hfill\break

\section{Productiviy and Value}\label{sec-productivityvalue}

Society's living requires:

\begin{itemize}
\tightlist
\item
  ``Resources'' are processed to create goods to be consumed
\item
  ``Labor'' is used to process goods and services
\item
  ``Capital'' intermediates value between goods and services
\end{itemize}

Society organizes in groups of households:

\begin{itemize}
\tightlist
\item
  natural households
\item
  juristic corporations
\item
  governmental institutions
\end{itemize}

Households cooperate though exchange (trade) of goods and services
between ``Demand'' and ``Supply'' using monetary valuation:

\begin{itemize}
\tightlist
\item
  Goods Market Production, Inventory, Trade \(\rightarrow\) ``Price''
  for value of production
\item
  Labor Market Employment, Demographic Structure, Work Capabilities
  \(\rightarrow\) ``Wage'' for value of labor
\item
  Capital Market Savings, Investments, Exchange \(\rightarrow\)
  ``Interest'' for value of capital
\end{itemize}

\hfill\break
acronyms see Section~\ref{sec-acroeconom}\\
\strut \\

Supply is provided by the income \(Y\):

\begin{equation}\phantomsection\label{eq-ecoincome}{\begin{split} &
Y = L  + Y_H + Y_A + T_A + Y_G - Z_G
\end{split}}\end{equation}

\begin{equation}\phantomsection\label{eq-gdpsupply}{\begin{split} &
Y = C + S + T + M - X
\end{split}}\end{equation}

Demand is satisfied by the production \(V_I\):

\begin{equation}\phantomsection\label{eq-gdpdemand}{\begin{split} &
V_I = C + I + G + X - M
\end{split}}\end{equation}

Equilibrium:

\begin{equation}\phantomsection\label{eq-ecoequi}{\begin{split} &
Y = V_I       \\&
L - C = S = I        \\&
C + S = L        \\&
C + I = R
\end{split}}\end{equation}

Domestic Production \& International Trade:

\begin{equation}\phantomsection\label{eq-gdp}{\begin{split} &
V_S = Y + T - G_A + D_A
\end{split}}\end{equation}

\begin{equation}\phantomsection\label{eq-nnp}{\begin{split} &
V_N = V_S - D_A
\end{split}}\end{equation}

\begin{equation}\phantomsection\label{eq-gdp}{\begin{split} &
V_I = V_S - R_M + R_X
\end{split}}\end{equation}

Monetary Value of Production:

\begin{equation}\phantomsection\label{eq-moneyquant}{\begin{split} &
P \cdot V_I = Q \cdot N
\end{split}}\end{equation}

Macro Productivity:
\begin{equation}\phantomsection\label{eq-productivity}{\begin{split} &
Productivity = \frac{Output}{Input} = \frac{V_I}{Y}
\end{split}}\end{equation}\\
\strut \\

We track production at microeconomic level through (acronyms see
Section~\ref{sec-acroeconom}):

\(A =\) Assets represent ``Property Capital'' which consist of ``Current
Capital'' and ``Structural Capital''. Those investments are required for
the production processes.

\(B =\) Liabilities represent ``Financial Capital'' which consist of
``Exogen Funds'' and ``Endogen Funds''. Later includes ``Equity
Capital'' and ``Retained Gains''. Those funds do finance Assets.

\(W =\) Expenses represent ``Work Efforts'' to ``Supply Production''
with use of Assets.

\(Y =\) Income represents ``Received Earnings'' achieved with production
to ``Satisfy Demand''. Parts of the gains will be retained to finance
Assets.

\(R =\) Gain, Profit or Loss from executed work (resultant value from
input efforts) is represented by \(R = Y - W\)

The ``transfer of energy'' between all subsystems is mapped in balancing
out theirs capital (usage of capital for production though possessed
assets): \(A = B + R\)

Micro Productivity:
\begin{equation}\phantomsection\label{eq-microprod}{\begin{split} &
Productivity = \frac{Output}{Input} = \frac{Y}{W}
\end{split}}\end{equation}

\hfill\break
\hfill\break

\section{Hypothesis and Conjecture}\label{sec-hypothesis}

Conjecture of Equation~\ref{eq-moneyquant} and
Equation~\ref{eq-workthermal}

How consistently map economic measures to the quantities of nature to
model the economical system?

Input:
\begin{equation}\phantomsection\label{eq-ecoinput}{\begin{split} &
Input = f_0(Resouces, Labor, Capital) \text{ = Work = } \Delta\text{Energy} = p \cdot v
\end{split}}\end{equation}

Output:
\begin{equation}\phantomsection\label{eq-ecooutput}{\begin{split} &
Output \text{ = Goods + Services = Production = } \Delta\text{State} = V_I 
\end{split}}\end{equation}

Productivity:
\begin{equation}\phantomsection\label{eq-ecooutput}{\begin{split} &
Productivity = \frac{Output}{Input} = \frac{\Delta State}{\Delta Energy} = \frac{Production}{Work} = \frac{V_I}{pv}
\end{split}}\end{equation}

see description of ``movement'' Equation~\ref{eq-movement}

\hfill\break
\hfill\break

\begin{longtable}[]{@{}
  >{\raggedright\arraybackslash}p{(\linewidth - 8\tabcolsep) * \real{0.1930}}
  >{\raggedright\arraybackslash}p{(\linewidth - 8\tabcolsep) * \real{0.2281}}
  >{\centering\arraybackslash}p{(\linewidth - 8\tabcolsep) * \real{0.1754}}
  >{\raggedleft\arraybackslash}p{(\linewidth - 8\tabcolsep) * \real{0.2105}}
  >{\raggedleft\arraybackslash}p{(\linewidth - 8\tabcolsep) * \real{0.1930}}@{}}
\caption{Core Measurements}\label{tbl-coremeasuremnts}\tabularnewline
\toprule\noalign{}
\begin{minipage}[b]{\linewidth}\raggedright
\textbf{Economy}
\end{minipage} & \begin{minipage}[b]{\linewidth}\raggedright
\(\rightarrow\)
\end{minipage} & \begin{minipage}[b]{\linewidth}\centering
\end{minipage} & \begin{minipage}[b]{\linewidth}\raggedleft
\(\leftarrow\)
\end{minipage} & \begin{minipage}[b]{\linewidth}\raggedleft
\textbf{Physics}
\end{minipage} \\
\midrule\noalign{}
\endfirsthead
\toprule\noalign{}
\begin{minipage}[b]{\linewidth}\raggedright
\textbf{Economy}
\end{minipage} & \begin{minipage}[b]{\linewidth}\raggedright
\(\rightarrow\)
\end{minipage} & \begin{minipage}[b]{\linewidth}\centering
\end{minipage} & \begin{minipage}[b]{\linewidth}\raggedleft
\(\leftarrow\)
\end{minipage} & \begin{minipage}[b]{\linewidth}\raggedleft
\textbf{Physics}
\end{minipage} \\
\midrule\noalign{}
\endhead
\bottomrule\noalign{}
\endlastfoot
Resources & Assets & & Force & Matter \\
Capital & Liabilities & & Momentum & Space \\
Labor & Expenses & & Work & \\
& Income & Production & Energy & Time \\
\end{longtable}

\hfill\break
\hfill\break

As a first approach, considering the geometrical perspective of the
concept of ``work'' (Equation~\ref{eq-workgeom}), applied to economical
measurements, so to satisfy the following system:

\(f_A(Assets) = f_B(Liabilities + Gain)\) = ``momentum'' \(p\) from
initial state \(r_0\) to new state \(r_1\), where the change of state is
\(r = r_1 - r_0\)

Matter substance, object movement/change:
\begin{equation}\phantomsection\label{eq-physassets}{\begin{split} &
A = f_A = F \cdot t = p
\end{split}}\end{equation}

Space of movement, maneuverability:
\begin{equation}\phantomsection\label{eq-physliabilities}{\begin{split} &
B + R = f_B = m \cdot v = p
\end{split}}\end{equation}

Work input (efforts, costs) through labor force, machines and
infrastructure applied, and materials consumed for production:
\begin{equation}\phantomsection\label{eq-physexpens}{\begin{split} &
W = f_0 = F \cdot r = m \cdot \frac{v}{t} \cdot r = p \cdot v
\end{split}}\end{equation}

Energy transfer, energy change, profit or loss:
\begin{equation}\phantomsection\label{eq-phyincome}{\begin{split} &
Y = \Delta E = \frac{1}{2} pv + \mathcal{U}
\end{split}}\end{equation}

Production value inflation adjusted:
\begin{equation}\phantomsection\label{eq-physprod}{\begin{split} &
P \cdot V_I = \sum Y_{adjust} = \sum (W + R)_{adjust} = \Delta U - \Delta H = Q \cdot N
\end{split}}\end{equation}

see Equation~\ref{eq-workthemostat} and Equation~\ref{eq-workgeom}

Thus, to be determined are \(v\) (economy velocity) and \(p\) (financial
structure). Eventually, determine \(r\) = ``state change through
trajectory'' (substance change to the new assets structure) by given
\(v\) = ``velocity of movement'' (Business speed = monetary value, price
with inflation), or determine \(v\) by given \(r\). The individual
positions \(\{ i,j,l,m \}\) of Liabilities \(\{ B_i + R_j \}_k\), Assets
\(\{ A_m \}_k\), and Expenses \(\{ W_l \}_k\) shall be grouped in sets
\(\{ k \}\), so to satisfy \(f_k(i,j) = f_k(m) = f_k(l) = \{ pv \}_k\).

The constraints shall be modelled with the Lagrangian. Hence, ``Work''
shall be derived from the Hamiltonian which leads to the identification
of the ``movement'' (= change, transformation) of the trajectory (= path
to new state) with \(\dot r = v\) and \(\dot p = F\).

Experiment design: \textbf{\(p\) and \(v\) are modelled ex-ante with
means of CONSTRAINTS from contractual agreement and from given juristic
legislation, in such a way that physics laws are satisfied. Decisions
making occure within this ex-ante frame.}

\hfill\break
\hfill\break

\section{Acronyms}\label{acronyms}

\subsection{\texorpdfstring{Acronyms of Physics:
Chapter~\ref{sec-workenergy}}{Acronyms of Physics: Chapter~}}\label{sec-acroworkenergy}

\begin{longtable}[]{@{}
  >{\raggedright\arraybackslash}p{(\linewidth - 2\tabcolsep) * \real{0.4510}}
  >{\raggedright\arraybackslash}p{(\linewidth - 2\tabcolsep) * \real{0.5490}}@{}}
\toprule\noalign{}
\endhead
\bottomrule\noalign{}
\endlastfoot
\(A =\) Action & \(E = E_{kin} + E_{pot} =\) total energy \\
\(t =\) time & \(r =\) dpace (distance between two points,
one-dimensional length) \\
\(v =\) velocity & \(c =\) light speed \\
\(p =\) momentum & \(F =\) force \\
\(I =\) inertia & \(\mathcal{L} =\) Lagrangian \\
\(\mathcal{H} =\) Hamiltonian & \(\mathcal{K} =\) kinetic Energy \\
\(\mathcal{U} =\) effective potential Energy (\(\in E_{pot}\)) &
\(\mathcal{V} =\) potential Energy (\(\in E_{pot}\)) \\
\(\mathcal{Z} = \frac{1}{2}\frac{L_0^2}{mr^2} =\) Centrifugal Potential
& \(V = r_1 \cdot r_2 \cdot r_3 =\) Volume \\
\(\mathcal{k} =\) Wave Vector (``curvature'') & \(W =\) Work (done
vs.~received) \\
\(P =\) Pressure & \(L_0 =\) angular momentum \\
\(T =\) endogen Temperature & \(H =\) exogen Heat \\
\(U =\) endogen Energy (\(E_{kin} + E_{pot}\)) &
\(\Phi = \frac{\mathcal{V}}{q} =\) Electric Potential \\
\(\mathcal{A} =\) Magnetic Potential & \(b_0 =\) Boltzmann constant \\
\(g_0 =\) Gas constant & \(m =\) mass \\
\(q =\) charge & \$\(n =\) amount of objectes (particles density
\(n=\frac{N}{V}\)), \(n=2\) and \(f=1\) \\
\({\epsilon}_0 =\) electric constant & \(f = 3n \pm z =\) degrees of
freedom, \(f=1\) and \(n=2\) \\
\({\mu}_0 =\) magnetic constant & \(\mathcal{B} =\) Magnetic Field \\
\(\mathcal{E} =\) Electric Field &
H\(\mathcal{C} = m \cdot c_0 = \frac{\Delta H}{\Delta T}  =\) heat
capacity \\
\(c_0 = \frac{1}{m}\frac{\Delta H}{\Delta T} =\) specific heat &
\(S = b_0 \cdot \ln(\Omega) =\) Entropy (macro state) \\
\(l =\) Moll quantity & \(\Omega =\) micro states \\
\(z =\) amount of constraints (boundry conditions) & \(b_0 =\) Boltzmann
constant \\
\(\kappa = \frac{c_P}{c_V} = \frac{f+2}{f} =\) adiabaty &
\(\iota = \kappa\) ↔ adiabat \\
\(\iota = 0\) ↔ isobar & \(\iota = 1\) ↔ isotherm \\
\(\iota = \infty\) ↔ isochor & \ldots{} \\
\ldots{} & \ldots{} \\
\end{longtable}

\subsection{\texorpdfstring{Acronyms of Economy:
Chapter~\ref{sec-productivityvalue}}{Acronyms of Economy: Chapter~}}\label{sec-acroeconom}

\begin{longtable}[]{@{}
  >{\raggedright\arraybackslash}p{(\linewidth - 2\tabcolsep) * \real{0.5000}}
  >{\raggedright\arraybackslash}p{(\linewidth - 2\tabcolsep) * \real{0.5000}}@{}}
\toprule\noalign{}
\endhead
\bottomrule\noalign{}
\endlastfoot
\(T =\) Taxes & \(M =\) Import of Goods and Services from foreing
symstes \\
\(G =\) Governent Expenses, incl.~Social Insurances & \(X =\) Export of
Goods and Services to foreign system \\
\(Y =\) Income of Economy from Turnover & \(G_A =\) Governental
Subsidies \\
\(D_A =\) Depreciations (Reinvestments) on Assets & \(V_N =\) Net
Naötional Production, Society NNP \\
\(N =\) Monetary Quantity & \(Q =\) Monetary Turnover Velocity \\
\(V_I =\) Gross Domestic Product GDP = \(\frac{Output}{Input}\),
Tradevolume & \(P =\) Price niveau (Inflation adjusted Value) \\
\(L =\) Wages from Labor Work (Salaries, \ldots) & \(R =\) Returns,
Earnings, Gains \\
\(Y_A =\) Income of priv. Business Households (Companies, Services, Real
Estate Rentals, Retained Profits) & \(Y_H =\) Income from priv. Capital
Households (Interests, Coupons, Dividends, \ldots{} of priv. Assets,
Investmens, Credits, Debits, Bonds, Equity) \\
\(T_A=\) Tax on Capital of Corprate Compaies (Business Assets) &
\(Y_G =\) Governmental Income from Assets, Services, Social
Institutions/Insurances \\
\(Z_G =\) Interests on Governmental Debt & \(V_S =\) Gross National
Produkt, Society GNP \\
\(I =\) Investments on Assets, incl.~Storage Change & \(R_M =\) Capital
Earnings and Wages from Abroad (from Foreign System) \\
\(R_X\) Capital Earnings and Wages to Abroad (to Foreign System) &
\(W =\) Expensens, costs from human and machinary work efforts \\
\end{longtable}

\hfill\break
\hfill\break

see also {[}1{]}

\begin{center}\rule{0.5\linewidth}{0.5pt}\end{center}

\phantomsection\label{refs}
\begin{CSLReferences}{0}{1}
\bibitem[\citeproctext]{ref-wachter_compendium_2005}
\CSLLeftMargin{1. }%
\CSLRightInline{Wachter A, Hoeber H (2005)
\href{https://link.springer.com/book/10.1007/0-387-29198-9}{Compendium
of {Theoretical} {Physics}}. Springer}

\bibitem[\citeproctext]{ref-cockshott_w_paul_classical_2011}
\CSLLeftMargin{2. }%
\CSLRightInline{Cockshott WP (2011)
\href{https://www.routledge.com/Classical-Econophysics/Cottrell-Cockshott-Michaelson-Wright-Yakovenko/p/book/9780415696463}{Classical
{Econophysics}}. Routledge Advances in Experimental; Computable
Economics}

\bibitem[\citeproctext]{ref-wladyslaw_macroeconometric_2013}
\CSLLeftMargin{3. }%
\CSLRightInline{Wladyslaw W (2013)
\href{https://link.springer.com/book/10.1007/978-3-642-34468-8}{Macroeconometric
{Models}}, 1st ed. Springer}

\bibitem[\citeproctext]{ref-lichtenegger_schlusselkonzepte_2015}
\CSLLeftMargin{4. }%
\CSLRightInline{Lichtenegger K (2015)
\href{https://link.springer.com/book/10.1007/978-3-8274-2385-6}{Schlüsselkonzepte
zur {Physik}}, 1st ed. Springer Spektrum}

\bibitem[\citeproctext]{ref-padmanabhan_sleeping_2015}
\CSLLeftMargin{5. }%
\CSLRightInline{Padmanabhan T (2015)
\href{https://link.springer.com/book/10.1007/978-3-319-13443-7}{Sleeping
{Beauties} in {Theoretical} {Physics}}, 1st ed. Springer}

\bibitem[\citeproctext]{ref-carter_foundations_2001}
\CSLLeftMargin{6. }%
\CSLRightInline{Carter M (2001)
\href{https://mitpress.mit.edu/9780262531924/foundations-of-mathematical-economics/}{Foundations
of {Mathematical} {Economics}}. The MIT Press}

\bibitem[\citeproctext]{ref-michaelides_21_2025}
\CSLLeftMargin{7. }%
\CSLRightInline{Michaelides PG (2025)
\href{https://link.springer.com/book/10.1007/978-3-031-76140-9}{21
{Equations} that {Shaped} the {World} {Economy}}, 1st ed. Palgrave
Macmillan Cham}

\bibitem[\citeproctext]{ref-pires_da_cruz_introduction_2025}
\CSLLeftMargin{8. }%
\CSLRightInline{Pires da Cruz J (2025)
\href{https://link.springer.com/book/10.1007/978-3-032-04346-7}{Introduction
to the {Physics} of the {Economy} and {Finance}}, 1st ed. Springer Cham}

\bibitem[\citeproctext]{ref-cass_hamiltonian_2014}
\CSLLeftMargin{9. }%
\CSLRightInline{Cass D, Shell K (2014)
\href{https://shop.elsevier.com/books/the-hamiltonian-approach-to-dynamic-economics/cass/978-0-12-163650-0}{The
{Hamiltonian} {Approach} to {Dynamic} {Economics}}, 1st ed. Elsevier,
Academic Press}

\bibitem[\citeproctext]{ref-haven_quantum_2017}
\CSLLeftMargin{10. }%
\CSLRightInline{Haven E, Khrennikov AI, Robinson T (2017)
\href{https://www.worldscientific.com/doi/epdf/10.1142/q0080}{Quantum
{Methods} in {Social} {Science} - {A} {First} {Course}}. World
Scientific, Publishing Europe Ltd}

\end{CSLReferences}



\printindex


\end{document}
